\chapter{CI/CD: Các khái niệm}
\label{ch:cicd}

\section{Hai thực hành, một pipeline}

\begin{description}
  \item[Continuous Integration] mọi thay đổi được merge vào branch
  chung thường xuyên và được \emph{kiểm chứng tự động} --- biên dịch,
  test, đóng gói. Sản phẩm của CI là sự tự tin: branch chính luôn ở
  trạng thái được biết là tốt, và một hỏng hóc bị bắt vài phút sau
  commit gây ra nó, khi tác giả vẫn còn nhớ ngữ cảnh.
  \item[Continuous Delivery/Deployment] mọi thay đổi đã kiểm chứng
  được tự động chuẩn bị để phát hành (Delivery) hoặc thật sự được
  phát hành (Deployment --- không có cổng chặn bởi con người). Dự án
  này thực hành dạng mạnh hơn: một lần push lên \code{main} kết thúc
  bằng các pod đang chạy, không ai bấm nút.
\end{description}

Pipeline là cỗ máy cưỡng chế cả hai:

\begin{figure}[h]
\centering
\begin{tikzpicture}[node distance=3.5mm and 5mm]
  \node[boxdark] (git) {git push};
  \node[box, right=of git] (trig) {trigger};
  \node[box, right=of trig] (build) {build\\\scriptsize + test};
  \node[box, right=of build] (pkg) {package\\\scriptsize image + tag};
  \node[box, right=of pkg] (reg) {registry};
  \node[box, right=of reg] (dep) {deploy\\\scriptsize + verify};
  \draw[flow] (git) -- (trig);
  \draw[flow] (trig) -- (build);
  \draw[flow] (build) -- (pkg);
  \draw[flow] (pkg) -- (reg);
  \draw[flow] (reg) -- (dep);
\end{tikzpicture}
\caption{Hình dạng pipeline phổ quát. Tên công cụ thay đổi; các stage thì không.}
\end{figure}

Dù là Jenkins+Jib+k3s (dự án này), GitHub Actions+Docker+EKS, hay
GitLab CI+Kaniko+OpenShift --- mã nguồn đi vào, artifact đã test đi
vào registry, bộ điều phối được bảo hội tụ về nó, rồi kiểm chứng rằng
nó đã làm vậy. Học hình dạng này một lần và pipeline của mọi nhà tuyển
dụng chỉ còn là bài tập đổi tên.

\section{Từ vựng}

\begin{description}
  \item[Trigger] thứ khởi động một lần chạy. \emph{Webhook}: máy chủ
  git gọi CI server khi có push --- tức thì, nhưng đòi hỏi CI server
  phải truy cập được từ internet. \emph{Polling}: CI hỏi repo ``có gì
  mới không?'' theo lịch --- không cần ai truy cập được, đổi lại là độ
  trễ. Dự án này dùng polling (\code{H/5 * * * *}: khoảng năm phút một
  lần) chính vì Jenkins nằm sau NAT ở nhà, không có địa chỉ công khai
  --- topology quyết định trigger, một mẫu hình bạn sẽ gặp ở mọi mạng
  doanh nghiệp bị khóa chặt.
  \item[Stage / step] các đơn vị của pipeline: stage là các pha có
  tên (hiện trên UI, chặn tiến trình); step là các lệnh bên trong.
  \item[Artifact] bất cứ thứ gì build sinh ra mà đáng giữ lại: file
  jar, image, báo cáo test. Artifact phải \emph{bất biến và truy vết
  được} --- vì thế image được tag bằng SHA của git, để một pod đang
  chạy có thể truy ngược về đúng commit mà nó đang chạy.
  \item[Test gate] tính chất rằng các stage sau chỉ chạy nếu các stage
  trước đã qua. Một pipeline triển khai dù test thất bại chỉ là một
  script triển khai với vài bước thừa.
  \item[Environment promotion] cùng một artifact đi từ dev
  $\rightarrow$ staging $\rightarrow$ prod. Quy tắc sắt: \emph{thăng
  cấp artifact, không bao giờ build lại} --- một lần build lại ``cho
  prod'' là một binary khác với cái bạn đã test. Bộ máy biến môi
  trường của Chương~1 tồn tại chính xác để một image vừa với mọi môi
  trường.
  \item[Rollback] lối thoát hiểm. Vì registry giữ các image tag theo
  SHA và Kubernetes giữ các ReplicaSet cũ (Chương~15), rollback là trỏ
  lại một artifact trước đó --- vài giây, không phải một lần build
  lại.
\end{description}

\begin{decision}[triển khai ở mọi lần push lên main --- cho một lập trình viên]
Continuous Deployment không cần phê duyệt là lựa chọn đúng ở đây: một
lập trình viên, bán kính ảnh hưởng cỡ homelab, và vòng phản hồi nhanh
nhất có thể. Cùng lựa chọn đó ở một ngân hàng sẽ sai, không phải vì tự
động hóa rủi ro mà vì các thay đổi \emph{chịu quản lý} cần bằng chứng
đã được xem xét. Con đường lên cấp giữ nguyên pipeline và thêm cổng:
một stage phê duyệt thủ công trước prod, hoặc --- Chương~27 --- GitOps,
nơi ``deploy'' trở thành một pull request có thể xem xét vào repo cấu
hình. Hình dạng pipeline sống sót; chỉ chính sách cổng chặn thay đổi.
\end{decision}

\begin{pitfall}[pipeline xanh mà nói dối]
Triệu chứng: lần chạy nào cũng xanh; production vẫn hỏng. Các nguyên
nhân kinh điển: test bị bỏ qua trong lệnh build (\code{-DskipTests}
--- điều mà pipeline \emph{này} làm ở stage build, chấp nhận một cách
có ý thức vì \code{mvn verify} chạy trong CI trên GitHub Actions ---
hãy biết rõ các cổng của chính mình); stage deploy không kiểm chứng
rollout (pipeline này có --- \code{rollout status} làm build thất bại
nếu pod không bao giờ Ready); hoặc thành công được đo bằng ``các lệnh
thoát với mã 0'' thay vì ``service trả lời request.'' Một pipeline
xanh là một lời khẳng định; hãy biết chính xác nó khẳng định điều gì.
\end{pitfall}

\begin{quiz}
\begin{enumerate}
  \item CI và CD đưa ra những lời hứa khác nhau. Nêu mỗi lời hứa
  trong một câu, và chỉ ra stage nào trong pipeline của dự án này
  thực hiện lời hứa nào.
  \item Webhook so với polling: nêu ràng buộc quyết định trong dự án
  này, và một bối cảnh doanh nghiệp có cùng ràng buộc đó.
  \item Vì sao một artifact được thăng cấp không bao giờ được build
  lại cho từng môi trường, và cơ chế nào từ Chương~1 làm cho việc
  thăng cấp một-artifact trở nên khả thi?
  \item Pipeline của bạn xanh và prod sập. Liệt kê ba cách khác nhau
  mà một pipeline có thể nói dối, và cách sửa cho mỗi cách.
\end{enumerate}
\end{quiz}
